\section{任意の pants 木から得られる分岐基底}\label{topo:section}

本節の結果は有限複体についての定理であり，シンプレクティックな接着の成立を仮定しない．
したがって，前節の大域構成とは独立に用いることができる．
以下，$d\geq5$，$[d]=\{1,\ldots,d\}$ とする．

\begin{definition}\label{topo:definition}
各非順序対 $ij=\{i,j\}\subset[d]$ に対し，葉が $[d]\setminus\{i,j\}$ でラベル付けされた
有限三価木 $T_{ij}=T_{ji}$ を選ぶ．内部頂点の集合を $P_{ij}$，内部辺の集合を $A_{ij}$ とし，
葉 $k$ に隣接する内部頂点を $p_{ij}(k)$ と書く．
分岐複体 $B(T)$ を次のセルで定める．
\begin{enumerate}
\item 頂点は $v_1,\ldots,v_d$ とする．
\item $p\in P_{ij}$ ごとに $v_i$ と $v_j$ を結ぶ辺 $e_{ij,p}$ を置く．
\item 内部辺 $pq\in A_{ij}$ ごとに，$e_{ij,p}$ と $e_{ij,q}$ に沿う二角形を一枚付ける．
\item $i<j<k$ ごとに，$e_{ij,p_{ij}(k)}$，$e_{jk,p_{jk}(i)}$，
$e_{ik,p_{ik}(j)}$ を境界とする三角形を一枚付ける．
\end{enumerate}
二角形は二本の異なる辺に沿う閉円板である．この記述は有限正則 CW 複体を定め，
細分すれば有限単体的多面体になる．
\end{definition}

三価木の辺数の公式から
\[
 |P_{ij}|=d-4,\qquad |A_{ij}|=d-5
\]
である．葉と葉に接する開辺を除く部分 $T_{ij}^{\rm int}$ は，内部頂点を全て含む木であり，
$d=5$ では一点である．

\begin{proposition}\label{topo:counts}
$B(T)$ の各辺の内部にはちょうど三枚の面が接する．セル数は
\[
 V=d,\qquad E=\binom d2(d-4),\qquad
 F_\triangle=\binom d3,\qquad F_2=\binom d2(d-5).
\]
頂点 $v_i$ のリンク $G_i$ は，完全グラフ $K_{d-1}$ の各頂点 $j$ を
$T_{ij}^{\rm int}$ で展開した連結三価グラフであり，
\[
 |V(G_i)|=(d-1)(d-4),\quad
 |E(G_i)|=\frac32(d-1)(d-4),\quad
 b_1(G_i)=\binom{d-2}{2}.
\]
\end{proposition}

\begin{proof}
$e_{ij,p}$ に接する面は，$T_{ij}$ における $p$ の三つの flags に対応する．
内部辺の flag は二角形，葉 $k$ の flag は三角形 $ijk$ を与えるため，面の数は三である．
リンクの頂点は $(j,p)$，$j\ne i$，$p\in P_{ij}$ である．
その辺は，各 $T_{ij}^{\rm int}$ の内部辺と，$j,k\ne i$ に対応する
$(j,p_{ij}(k))$ と $(k,p_{ik}(j))$ を結ぶ辺である．
内部木を全て縮約すると $K_{d-1}$ になり，連結性と第一 Betti 数が従う．
三価性から $2|E(G_i)|=3|V(G_i)|$ も従う．
\end{proof}

\subsection{縮約写像と整数係数の鎖収縮}

$S_d=\operatorname{Sk}_2\Delta^{d-1}$ と置く．同じ $ij$ の平行辺を $S_d$ の辺 $ij$ に送り，
二角形をその辺に潰し，三角形を対応する三角形に写す写像
\[
 q:B(T)\longrightarrow S_d
\]
がある．より具体的には，$ij$ の全ての辺と二角形の和集合を
\[
 \bigl(T_{ij}^{\rm int}\times[0,1]\bigr)/
 \bigl(T_{ij}^{\rm int}\times\{0\}\sim v_i,\ 
 T_{ij}^{\rm int}\times\{1\}\sim v_j\bigr)
\]
と表示し，第二因子への射影を使う．三角形上の写像は，この境界のパラメータ表示と整合するように選ぶ．
各辺の内点の逆像は $T_{ij}^{\rm int}$，それ以外の点の逆像は一点である．
したがって $q$ は cell-like な写像，すなわち全ての点逆像が可縮な compact 多面体である写像である．

\begin{theorem}\label{topo:sdr}
任意の木の組 $T=(T_{ij})$ に対し，$q$ はホモトピー同値であり，
\[
 B(T)\simeq\bigvee^{\binom{d-1}{3}}S^2.
\]
さらに，セル鎖写像 $q_*$ は整数係数で明示的な鎖切断と鎖ホモトピーを持つ．
\end{theorem}

\begin{proof}
$i<j$ に対して $e_{ij,p}$ を $v_i$ から $v_j$ に向ける．各内部木に根
$r_{ij}\in P_{ij}$ を選び，内部辺を根から遠ざかる向きに向ける．
その向きが $p\to u$ である二角形を $b_{ij,pu}$ と書き，
\[
 \partial b_{ij,pu}=e_{ij,u}-e_{ij,p}
\]
となる向きを与える．根から $p$ への経路に沿う二角形鎖を $D_{ij}(p)$ とすると
\[
 D_{ij}(r_{ij})=0,\qquad
 \partial D_{ij}(p)=e_{ij,p}-e_{ij,r_{ij}}.
\]
三角形 $ijk$ の向きを $i<j<k$ の順序で定め，その鎖を $\tau_{ijk}$ とする．
補正三角形
\begin{align*}
 \widehat\tau_{ijk}={}&\tau_{ijk}
   -D_{ij}(p_{ij}(k))-D_{jk}(p_{jk}(i))
   +D_{ik}(p_{ik}(j))
\end{align*}
は
\[
 \partial\widehat\tau_{ijk}
 =e_{ij,r_{ij}}+e_{jk,r_{jk}}-e_{ik,r_{ik}}
\]
を満たす．したがって，$S_d$ の頂点を同じ頂点に，辺 $ij$ を
$e_{ij,r_{ij}}$ に，三角形 $ijk$ を $\widehat\tau_{ijk}$ に送る写像
$s:C_*(S_d;\ZZ)\to C_*(B(T);\ZZ)$ は鎖写像であり，$q_*s=1$ である．

次数一の生成元に対し $H(e_{ij,p})=D_{ij}(p)$ と定め，それ以外の次数で
$H=0$ と置く．木の経路の一意性から，内部辺 $p\to u$ に対して
$D_{ij}(u)-D_{ij}(p)=b_{ij,pu}$ である．これと補正三角形の式を使うと，
全てのセルについて
\[
 1-sq_* =\partial H+H\partial
\]
が成り立つ．また $Hs=0$，$q_*H=0$，$H^2=0$ である．
これは整数係数の鎖複体の明示的な strong deformation retraction である．

基本群について，二角形の関係は同じ $ij$ の全ての辺を相対端点ホモトピーで同一視する．
残る三角形の関係は完全グラフの全てのループを消すので，$B(T)$ は単連結である．
$S_d$ も単連結であるため，いま証明した整数ホモロジー同値と
単連結 CW 複体に対するホモロジー版 Whitehead 定理から，$q$ はホモトピー同値である
\cite[Corollary 4.33]{Hatcher}．
最後に，$S_d$ は単連結な二次元複体であり，$H_2(S_d;\ZZ)$ は自由群である．
その基底を Hurewicz の定理で球面からの写像として実現し，再び Whitehead の定理を使えば，
主張した球面の wedge を得る．球面数は
$\chi(S_d)-1=d-\binom d2+\binom d3-1=\binom{d-1}{3}$ である．
\end{proof}

\begin{corollary}[整数二サイクルの基底]\label{topo:cycles}
$2\leq i<j<k\leq d$ に対し
\[
 Z_{ijk}=\widehat\tau_{ijk}-\widehat\tau_{1jk}
              +\widehat\tau_{1ik}-\widehat\tau_{1ij}
\]
と置くと，これらは $H_2(B(T);\ZZ)$ の基底を与える．
\end{corollary}

\begin{proof}
標準複体では，$[ijk]-[1jk]+[1ik]-[1ij]$ は四面体 $[1ijk]$ の境界である．
任意の二サイクルから，頂点 $1$ を含まない各三角形の係数だけこれらを差し引くと，
残りは $1$ を含む三角形だけに支持される．そのような二サイクルは零である．
実際，辺 $ij$，$i,j>1$ はその支持の中では三角形 $1ij$ にだけ現れる．
同じ理由で基底の一次独立性も分かる．鎖切断 $s$ で移せばよい．
\end{proof}

この証明は有限個の次数の行列の階数計算を一般次数の証明に代用しない．
各木について根と経路を列挙するだけで，整数係数の切断・ホモトピー・二サイクルを出力できる．

\subsection{通常の双対複体との差}

\begin{proposition}\label{topo:nonhomeomorphic}
$d=5$ では $q:B(T)\to S_5$ は同相写像である．$d\geq6$ では
$B(T)$ と $S_d$ は同相でない．
\end{proposition}

\begin{proof}
$d=5$ では各内部木は一点であり，二角形が存在しない．
$d\geq6$ では点 $x$ の局所ホモロジー
$H_2(Y,Y\setminus\{x\};\ZZ)$ の階数で区別できる．
三枚の面が接する $B(T)$ の辺内ではこの階数は二である．一方 $S_d$ では，
面内・辺内・頂点における階数はそれぞれ
\[
 1,\qquad d-3,\qquad \binom{d-2}{2}
\]
である．これは点のリンクの第一 Betti 数から従う．$d\geq6$ ではどれも二でない．
\end{proof}

したがって，木による三価化は通常の細分ではない．
自由な辺がないこと，双対複体とホモトピー同値であること，双対複体と同相であることは
区別しなければならない．また，この議論自体は integral-affine atlas を構成しない．

\begin{corollary}[KE 退化の cubical candidate への位相的射]
\label{topo:KE-cubical-comparison}
$d\ge5$ に対し，別の研究で用いる有限立方体複体
\[
 \mathcal C_d=
 \left\{w\in[0,1]^d:\max_i w_i=1,\ 
             |\operatorname{supp}(w)|\le3\right\}
\]
を標準 square の貼合せで定める．正規化
\[
 \rho_d:\mathcal C_d\longrightarrow S_d,\qquad
 w\longmapsto\frac{w}{\sum_iw_i}
\]
は同相であり，
\[
 p_T=\rho_d^{-1}\circ q:B(T)\longrightarrow\mathcal C_d
\]
は全ての点逆像が可縮な連続全射で，整数ホモロジーの同型を誘導する
ホモトピー同値である．$d=5$ では $p_T$ は同相，
$d\ge6$ では同相でない．
\end{corollary}
\begin{proof}
$\sum_iw_i\ge1$ であり，$S_d$ の各点 $x$ は
$x/\max_i x_i\in\mathcal C_d$ を一意に定めるので $\rho_d$ は同相である．
したがって $p_T$ の点逆像は $q$ の点逆像と一致し，
定理\ref{topo:sdr} の鎖ホモトピーとホモトピー同値がそのまま適用できる．
同相性の区別は命題\ref{topo:nonhomeomorphic}から従う．
\end{proof}

\begin{lemma}[連結 fiber を持つ閉写像の合成]
\label{topo:connected-preimages}
$X$ を compact Hausdorff 空間，$Y$ を Hausdorff 空間とし，
$p:X\to Y$ を連結 fiber を持つ連続全射とする．
任意の連結部分集合 $A\subset Y$ に対し，$p^{-1}(A)$ は連結である．
従って compact Hausdorff 空間の間の連結 fiber の連続全射を合成しても，
連結 fiber の連続全射となる．
\end{lemma}
\begin{proof}
$p$ は閉写像である．$p^{-1}(A)$ 内の相対閉集合は
$C\cap p^{-1}(A)$，$C\subset X$ 閉，と書け，その像は $p(C)\cap A$ である．
従って制限 $p_A:p^{-1}(A)\to A$ も閉写像である．
$p^{-1}(A)=U\sqcup V$ が非自明な分離なら，連結な各 fiber は
$U,V$ の一方に全体として入る．よって $p(U),p(V)$ は
$A$ を覆う非空な互いに素な相対閉集合となり，連結性に反する．
合成の点逆像は，後の写像の連結な点逆像を前の写像で引き戻したものなので，
最初の主張を適用すればよい．
\end{proof}

\begin{corollary}[Cubical 射の全 fiber の連結性]
\label{topo:cubical-composite-connected}
本稿で構成する $\pi:X\to B_d(T)$ に対し，$f=p_T\pi:X\to\mathcal C_d$ は
proper な連結 fiber の連続全射である．特に
\[
 f^{-1}(w)=\pi^{-1}\bigl(p_T^{-1}(w)\bigr)
\]
は，$p_T^{-1}(w)$ が内部木である場合も連結である．
\end{corollary}
\begin{proof}
$\pi$ の主張は補題\ref{glue:closed-whole-fiber-quotient}の適用，
$p_T$ の主張は系\ref{topo:KE-cubical-comparison}による．
前補題を用い，properness は compactness から従う．
\end{proof}

$p_T$ の点逆像が可縮であることから $f$ の点逆像の可縮性は結論しない．
$d\ge6$ で内部木が辺を持つ場合，その開辺上の正則 torus の合併は
局所的に interval と $T^2$ の積であり，$f$ の fiber は三次元部分を含み得る．
Lagrangian torus fibration と呼ぶのは $\pi$ であり，
$f$ はその全 fiber 合併を扱う計量比較写像である．

\begin{remark}[比較の射程]
系\ref{topo:KE-cubical-comparison}は候補極限複体と Lagrangian 分岐基底の
\emph{labelled topology} の比較である．
$q$ は一般に距離を保存せず，$p_T$ から
KE 計量の GH 収束，action-compatible integral-affine atlas，
あるいは special Lagrangian 構造は導けない．
$d\ge6$ の非同相性は，両基底をそのまま同じ位相空間として
同一視することに対する具体的な障害である．
\end{remark}


\paragraph{既知の cubical subdivision との関係．}
\label{topo:cubical-barycentric-comparison}
$\mathcal C_d$ の各閉 cell は，空でない集合の組
$A\subseteq B\subseteq\{1,\ldots,d\}$，$|B|\le3$ に対して
\[
 Q_{A,B}=\{w_i=1\ (i\in A),\quad 0\le w_i\le1\ (i\in B\setminus A),
 \quad w_i=0\ (i\notin B)\}
\]
と書ける．その face poset は，単体の二骨格の非空 face poset の
閉区間 $[A,B]$ を包含で順序付けたものに等しい．
従って組合せ型は既知の cubical barycentric subdivision，すなわち
Babson--Billera--Chan\cite[1996年preprint, \S2.3, Proposition~2.2, Corollary~2.4]{BabsonBilleraChan}
の barycentric cover の特殊例である．
この組合せ的構成自体の新規性を主張しない．以下では各二次元 cell を
\emph{unit Euclidean square} とする長さ距離を固定し，
その非伸長 retraction と直径を別に証明する．距離結果の先行研究との
網羅的な照合は未了である．

\subsection{座標部分複体の非伸長 retraction}

以後，各 unit square の Euclidean 長さを貼り合わせた距離を $d_\square$ とし，
$Q_{p;qr}=e_p+[0,1]e_q+[0,1]e_r$ と書く．
非空な $I\subset\{1,\ldots,d\}$ に対して
\[
 \mathcal C_I=\{w\in\mathcal C_d:\operatorname{supp}(w)\subset I\}
\]
とし，その cubical cells の Euclidean 長さから定まる内在距離を $d_I$ と書く．
$|I|=1$ の場合は一点，$|I|=2$ の場合は二本の unit intervals の和である．
$|I|\ge3$ の場合は各 unit square を貼り合わせた有限長さ空間となる．
この定義だけでは $d_I$ は $d_\square$ の制限を意味しないが，
以下の直接構成によって両者は等しい．

\begin{proposition}[座標部分複体の非伸長 retraction]
\label{topo:coordinate-retraction}
$a\in I$ を一つ選び，$r_{I,a}:\mathcal C_d\to\mathbb R^d$ を
\[
 \bigl(r_{I,a}(w)\bigr)_b=
 \begin{cases}
  \min\!\left\{1,
   \left(w_a^2+\displaystyle\sum_{\ell\notin I}w_\ell^2\right)^{1/2}
       \right\},&b=a,\\
  w_b,&b\in I\setminus\{a\},\\
  0,&b\notin I
 \end{cases}
\]
で定める．これは $(\mathcal C_d,d_\square)$ から
$(\mathcal C_I,d_I)$ への $1$-Lipschitz retraction である．特に
\begin{equation}\label{topo:eq:coordinate-distance}
 d_I(x,y)=d_\square(x,y)\qquad(x,y\in\mathcal C_I).
\end{equation}
さらに，各 unit square $Q\subset\mathcal C_d$ 自身も
非伸長 retraction の像であり，$Q$ 上で $d_\square$ は
その Euclidean 距離に等しい．
\end{proposition}
\begin{proof}
外側の全座標と $a$ 座標を一つの非負座標にまとめるので，
非零座標の個数は増えず，像の support は $I$ に含まれる．
入力の値 $1$ を取る座標が $I\setminus\{a\}$ にあればそのまま残り，
それ以外なら像の $a$ 座標が $1$ となる．
したがって像は $\mathcal C_I$ に入る．
$w\in\mathcal C_I$ では $w_a\le1$ より $r_{I,a}(w)=w$ である．

写像 $u\mapsto\min\{1,\|u\|_2\}$ は Euclidean 距離に関して
$1$-Lipschitz である．これは逆三角不等式と，実数の clipping が
$1$-Lipschitz であることから従い，$u=0$ と $\|u\|_2=1$ でも成立する．
その微分が存在する点では勾配の長さは $1$ 以下であり，
clipping の外側では零である．変更しない他の座標はこの $u$ と直交するから
\[
 \|r_{I,a}(w)-r_{I,a}(w')\|_2^2
 \le\sum_{b\in I\setminus\{a\}}|w_b-w'_b|^2
       +|w_a-w'_a|^2+\sum_{\ell\notin I}|w_\ell-w'_\ell|^2
 =\|w-w'\|_2^2.
\]

ここでは target の内在距離も確認する必要がある．
一つの source square の固定座標を $p$ とする．
$p\in I\setminus\{a\}$ なら像でも $p$ 座標が $1$ であり，
$p=a$ または $p\notin I$ なら像では $a$ 座標が $1$ である．
残る可変座標は高々二つで，その label は元の square と $I,a$ から固定される．
従って各 source square は一つの target cell（square，辺，または点）に写る．
上の Euclidean 評価から，その square 内の任意の曲線の長さは増えない．
有限個の cell 内線分をつないだ経路に適用して長さを足し，
piecewise Euclidean 距離の定義に従って下限を取れば
\[
 d_I(r_{I,a}(x),r_{I,a}(y))\le d_\square(x,y)
\]
を得る．従って任意の rectifiable curve についても，距離空間の
$1$-Lipschitz 写像の一般的性質として長さは増加しない．
包含写像による逆向きの不等式と retraction 性から
\eqref{topo:eq:coordinate-distance} が従う．

最後に，$Q=\{w_p=1,\ 0\le w_q,w_r\le1,\ w_b=0
\text{ for }b\notin\{p,q,r\}\}$ に対して
\[
 P_Q(w)=e_p+w_qe_q+w_re_r
\]
を取る．これは $Q$ 上恒等な Euclidean 非伸長写像で，
target の $Q$ 自身は凸な Euclidean square なので
$P_Q:(\mathcal C_d,d_\square)\to Q$ も非伸長である．
従って $Q$ 上の全体距離と Euclidean 距離は等しい．
\end{proof}

\begin{corollary}[六つ以下の label による最短路]
\label{topo:six-label-gallery}
$d\ge3$ とする．任意の二点 $x,y\in\mathcal C_d$ は，
高々六つの label を持つ座標部分複体内の最短路で結べる．
その最短路は
\[
 N_d^*=3\binom{\min\{d,6\}}{3}
\]
本以下の unit-square 内線分からなる gallery として取れる．
特に $d=5$ では $N_d^*=30$，$d\ge6$ では $N_d^*=60$ である．
また $D_d=\operatorname{diam}(\mathcal C_d,d_\square)$ とすれば
$D_d=D_6$ for $d\ge6$．
\end{corollary}
\begin{proof}
$\operatorname{supp}(x)\cup\operatorname{supp}(y)$ は高々六元である．
これを含み，$3\le|I|\le\min\{d,6\}$ となる $I$ を選ぶ．
有限連結 cubical length space $\mathcal C_I$ は compact であり，最短路を持つ．
命題\ref{topo:coordinate-retraction}によって，その最短路は全体でも最短である．

全体の最短路が同じ unit square の二点を通れば，その間の長さは
両端間の Euclidean 距離に等しい．Euclidean 空間の三角不等式の等号条件より，
その区間は square 内の線分である．従って各 square の逆像は閉区間であり，
同じ square を出て再訪することはない．有限個のこれらの区間を使って
parameter interval を順に覆えば，各 square を高々一度用いる gallery を得る．
$\mathcal C_I$ の square の個数は $|I|\binom{|I|-1}{2}=3\binom{|I|}{3}$
なので所望の上界を得る．

$d\ge6$ のとき，任意の二点間距離は label の置換により
$\mathcal C_6$ 内の二点間距離に等しいから $D_d\le D_6$．
逆向きは $\mathcal C_6\subset\mathcal C_d$ の等長性から従う．
\end{proof}


\subsection{成分頂点の計量錐と固定補空間の無限遠}

\begin{proposition}[成分頂点近傍の厳密な錐模型]
\label{topo:component-cone-germ}
$i\in[d]$ を固定し，
\[
 C_{2,d-1}=\{u\in\mathbb R_{\ge0}^{[d]\setminus\{i\}}:
                   |\operatorname{supp}(u)|\le2\}
\]
に各座標 quadrant の Euclidean 長さを貼った距離 $d_C$ を入れる．
任意の $0<r<1/2$ に対し，
\[
 \iota_i(u)=e_i+\sum_{k\ne i}u_ke_k
\]
は閉球 $\overline B_{d_C}(0,r)$ から
$\overline B_{d_\square}(e_i,r)$ への等長全単射である．
両閉球の距離は，それぞれ錐全体と $\mathcal C_d$ 全体の距離の制限とする．
従って $e_i$ における一意な metric tangent cone は $C_{2,d-1}$ であり，
その link は各辺の長さが $\pi/2$ の完全グラフ $K_{d-1}$ である．
\end{proposition}
\begin{proof}
$\mathcal C_d$ の Euclidean 座標表示は曲線の長さを増加させない．
$w_i<1$ なら別の座標 $w_k=1$ が存在するため，
$d_\square(e_i,w)\ge\|e_i-w\|_2\ge1$ である．
一方，閉 star $S_i=\{w\in\mathcal C_d:w_i=1\}$ 内では，
$u=w-e_i$ に沿う動径線分と Euclidean 下界により
$d_\square(e_i,w)=\|u\|_2$ となる．錐側も $d_C(0,u)=\|u\|_2$ なので，
両閉球は集合として上の写像で一致する．

錐の半径 $r$ への動径射影は，各 quadrant では Euclidean 閉球への射影であり，
共通の face 上で一致する．従って piecewise Euclidean 長さに関して非伸長である．
二端点の norm が $r$ 以下なら，錐内の距離の下限は同じ閉球内の経路で取れ，
その値は原点経由の $2r$ 以下である．これらの経路を $S_i$ へ写せるので，
$d_\square(\iota_i(u),\iota_i(v))\le d_C(u,v)$ である．

逆に $S_i$ 内の経路は同じ長さの錐内経路になる．$S_i$ の外へ出る経路は，
退出時と再入時に，$w_i=w_k=1$ を満たすいずれかの gate を通る．
二端点の各可変座標は $r$ 以下だから，gate までの片側の長さは少なくとも $1-r$ である．
従ってこの経路の全長は $2(1-r)>2r\ge d_C(u,v)$ であり，短縮を与えない．
全経路で下限を取り逆向きの不等式を得る．
閉球の等長は任意に小さい $r$ で成り立つため tangent cone の主張も従う．
座標軸が link の頂点，二軸の quadrant が長さ $\pi/2$ の辺を与える．
\end{proof}

全平面退化の成分を $H_i=\{L_i=0\}\subset\PP^3$ と書くと，
\[
 U_i=H_i\setminus\bigcup_{k\ne i}(H_i\cap H_k),\qquad
 K_{H_i}+\sum_{k\ne i}(H_i\cap H_k)=(d-4)H|_{H_i}.
\]
すなわち $U_i$ は $\PP^2$ から一般位置の $d-1$ 本の直線を除いた固定補空間である．
固定空間の漸近錐を扱う companion 原稿
\cite[\texttt{cor:arrangement-cone}]{MultiscaleKECurrent}の $(n,m)=(2,d-1)$ は，
同じ $\operatorname{Ric}=-1$ の規格化で上の $C_{2,d-1}$ を与える．
軸の label は二重線 $H_i\cap H_k$，quadrant の label は三重点 $H_i\cap H_k\cap H_l$ に対応する．
従って固定成分の無限遠模型と compact な cubical 極限の成分頂点の接錐は，
計量と labels を保って一致する．この有限複体の命題だけから，
actual 退化族の任意の moving 領域の tensor 比較や中間尺度の収束は従わない．
後者について，本稿の大域 distortion を別に使う帰結を
系\ref{ke:intermediate-component-cone}で述べる．

\subsection{全ての局所錐と移動基点の拡大極限}

原点を共有する $m$ 本の半直線を，その path metric で貼った空間を
$\mathcal Y_m$ と書く．また $\mathcal K_{3\pi/2}$ は，三つの Euclidean quadrants
を循環的に貼った全角 $3\pi/2$ の Euclidean cone とする．
以下の積には Euclidean product distance を入れる．

\begin{proposition}[Cubical target の四種類の計量的局所模型]
\label{topo:cubical-local-models}
任意の $q\in\mathcal C_d$ について，十分小さい半径の閉球は，次の空間の
指定された中心の閉球と厳密に等長である．距離はいずれも全空間の ambient 距離の制限である．
\begin{enumerate}
\item $q=e_i$ のときは $(C_{2,d-1},0)$．
\item $|\operatorname{supp}q|=2$ のときは
      $(\mathbb R\times\mathcal Y_{d-2},(0,0))$．
      これには $q=e_i+e_j$ も含む．
\item $q=e_i+e_j+e_k$ のときは $(\mathcal K_{3\pi/2},0)$．
\item 残る点，すなわち support が三元で，その三座標が全て $1$ ではない点では
      $(\mathbb R^2,0)$．
\end{enumerate}
従ってこれらが各点の一意な metric tangent cone である．
\end{proposition}
\begin{proof}
まず局所的な貼合せが ambient 距離も決めることを確認する．有限個の closed squares
のうち $q$ を含まないもの，および incident square の $q$ を含まない faces は，
$q$ から正の Euclidean 距離だけ離れている．その全てを避ける $a>0$ を取ると，
$q$ の Euclidean $a$ 近傍は incident tangent sectors の $a$ 近傍と，
各 sector の距離と貼合せを保って同一視できる．Sector 上の半径 $r$ への動径射影は
Euclidean 閉球への射影であり，共通 faces 上で一致するため長さを増加させない．
よって $r<a/2$ の二端点は，局所錐の半径 $r$ の球内で長さ $2r$ 以下の経路で結べる．
一方，$a$ 近傍の外を通る経路の長さは $2(a-r)>2r$ 以上である．
従って小球上の全点対距離には近傍外からの短縮がなく，局所錐との厳密な等長を得る．

成分頂点では命題\ref{topo:component-cone-germ}を用いる．Support が $\{i,j\}$ の点は
長さ $2$ の折れ辺 $E_{ij}(u)$，$0<u<2$ 上にある．ここで
\[
 E_{ij}(u)=\begin{cases}e_i+u e_j,&0\le u\le1,\\
                      (2-u)e_i+e_j,&1\le u\le2.\end{cases}
\]
各 $k\notin\{i,j\}$ について $E_{ij}(u)+v e_k$，$v\ge0$ が一枚の局所半平面を与える．
$u=1$ では二つの quadrants が直線境界で貼られ，一つの半平面になる．
全ての $d-2$ 枚が $v=0$ の同じ直線だけを共有するので，第二の積模型を得る．

三座標が全て $1$ の点では，その一つを $1$ に保ち残り二つを減らす三つの
quadrants が現れる．三本の負向き座標軸を境界として循環的に貼られるので，
link は長さ $3\pi/2$ の circle である．Support が三元の残りの点では，
$1$ の座標が一つなら一枚の square の内部であり，二つなら二枚の half-planes の
共通境界にある．いずれも局所錐は Euclidean plane となる．
\end{proof}

\begin{proposition}[移動基点の全ての拡大極限]
\label{topo:cubical-moving-blowups}
$\lambda_\nu\to\infty$ と任意の $q_\nu\in\mathcal C_d$ を取る．
Pointed metric spaces $(\mathcal C_d,\lambda_\nu d_\square,q_\nu)$ は
部分列を取れば pointed GH 収束し，その全ての可能な極限は，等長の違いを除いて
\[
 (C_{2,d-1},z),\qquad
 (\mathbb R\times\mathcal Y_{d-2},(0,z)),\qquad
 (\mathcal K_{3\pi/2},z),\qquad
 (\mathbb R^2,0)
\]
である．各 $z$ は対応する錐又は $\mathcal Y_{d-2}$ の任意の有限の点を取る．
任意に指定した $\lambda_\nu\to\infty$ に対し，これらの各基点付き空間を
実現する $q_\nu$ を選べる．
\end{proposition}
\begin{proof}
Compactness により $q_\nu\to q$ としてよい．前命題の $q$ における固定した
局所錐 $Z$ で $q_\nu$ を表し，その座標を $\lambda_\nu$ 倍した点を $z_\nu$ とする．
前命題の full-ambient 等長を半径 $\rho>0$ の閉球上で固定する．任意の固定 $A$ に対し，
$d_Z(0,z_\nu)/\lambda_\nu\to0$ なので，十分大きい $\nu$ では
$d_Z(0,z_\nu)+A<\lambda_\nu\rho$ である．従って二つの中心を動かした半径 $A$ の球は，
この等長を rescale した写像の定義域と像に含まれる．全点対距離も既に保存されているから，
その閉球は $(Z,z_\nu)$ の閉球と厳密に等長である．
以下ではこの四つの $Z$ の移動基点だけを調べる．

$Z=\mathbb R^2$ では平行移動を使う．
$Z=\mathbb R\times\mathcal Y_{d-2}$ では第一因子の平行移動を除き，
第二因子の半直線を部分列で固定する．その座標が有限の値へ収束すれば同じ積の
その点を基点とする極限を得る．無限大へ向かえば，各固定球は半直線の頂点を避け，
二次元 Euclidean 球となる．

$Z=\mathcal K_{3\pi/2}$ では，$z_\nu$ が有界なら properness により収束部分列を取れる．
原点からの距離が無限大へ向かえば，各固定球は Euclidean 球になる．実際，
単位半径の circle は compact で，全ての点が平坦な近傍を持つため，共通の
正の flat-ball radius がある．錐の dilation によってこの半径は基点の半径に比例する．

$Z=C_{2,d-1}$ でも有界な $z_\nu$ は有限の基点を与える．それ以外では support を
部分列で固定し，さらに各非零座標が有限値へ収束するか無限大へ向かうよう選ぶ．
正座標は高々二つである．二座標が無限大へ向かう場合は，固定球の局所模型は
$\mathbb R^2$ である．一座標だけが無限大へ向かう場合は，その座標を平行移動し，
残りの座標は高々一つだけ正であることから
$\mathbb R\times\mathcal Y_{d-2}$ を得る．
ここで単に closed star の内在距離を使っているのではない．発散する各座標の
最小値を $m_\nu$ とすると，半径 $A$ の二端点を結び，その正座標の star を
出る経路は長さ $2(m_\nu-A)>2A$ を要する．Star 内の積座標による閉球の経路は
長さ $2A$ 以下であり，外部の経路は短縮を与えない．これは固定配置錐の
positive-star 評価\cite[\texttt{lem:star-ball-isometry}]{MultiscaleKECurrent}と同じ機構である．

最後に各型を実現する．対応する前命題の点 $q$ を一つ固定し，その局所模型内で
指定した有限基点を $1/\lambda_\nu$ 倍した点を $q_\nu$ とする．
平坦な型は任意の一枚の square の内点を固定すればよい．固定球の等長により
所要の pointed 収束を得る．
\end{proof}

\subsection{六次以上の cubical target の厳密直径}

\begin{lemma}[三 label corner の singleton からの距離]
\label{topo:corner-singleton-radius}
$J=\{i,j,k\}$ とし，$x\in\mathcal C_J$ とする．すると
\[
 d_J(e_i,x)\le\sqrt5.
\]
等号は $x=e_j+e_k$ の場合に限る．特に $x_i>0$ なら不等号は真である．
\end{lemma}
\begin{proof}
$x_i=1$ なら，同じ square 内の線分を使って距離は $\sqrt2$ 以下である．
それ以外では $j,k$ を交換して $x_j=1$ とできる．
$u=x_i$，$v=x_k$ と置く．$w_i=1$ の square と $w_j=1$ の square を
共有辺で開いて $2\times1$ rectangle にすると，両点の差は $(2-u,v)$ となる．
この rectangle 内の線分を戻して
\[
 d_J(e_i,x)^2\le(2-u)^2+v^2\le5
\]
を得る．最後の等号には $u=0$，$v=1$ が必要である．

$x=e_j+e_k$ の下界も直接確認できる．各 square を singleton--triple
対角線で二分し，頂点値
\[
\begin{array}{c|ccccc}
z&e_i&e_j,e_k&e_i+e_j,e_i+e_k&e_j+e_k&e_i+e_j+e_k\\\hline
\sqrt5\,h(z)&0&4&2&5&3
\end{array}
\]
を affine に延長する．各三角形の二つの直交辺方向の増分は，
符号を除いて $1/\sqrt5$ と $2/\sqrt5$ である．
従って $h$ は $1$-Lipschitz であり，
$h(e_j+e_k)-h(e_i)=5/\sqrt5=\sqrt5$ が所望の下界を与える．
\end{proof}

\begin{theorem}[六 label からの直径の安定値]
\label{topo:diameter-six-labels}
任意の $d\ge6$ に対し
\[
 \operatorname{diam}(\mathcal C_d,d_\square)=2\sqrt5=\sqrt{20}.
\]
距離が $2\sqrt5$ に等しい点対は，互いに素な三元集合 $A,B$ による
\[
 v_A=\sum_{i\in A}e_i,\qquad v_B=\sum_{i\in B}e_i
\]
の対に限る．従ってこの直径と等号点対の記述は，KE の解析的な仮定と独立である．
\end{theorem}
\begin{proof}
まず $S=\operatorname{supp}(x)$，$T=\operatorname{supp}(y)$ とし，
$|S\cup T|\le5$ の場合を考える．
$S\cap T\ne\varnothing$ なら $i\in S\cap T$ を選び，
$S,T$ をそれぞれ含む三元集合 $J,K$ を取る．
$S\cap T=\varnothing$ なら少なくとも一方の大きさは二以下なので，
必要なら $x,y$ を交換して $|S|\le2$ とする．
$i\in T$ を選び，$J\supset S\cup\{i\}$，$K\supset T$ を三元集合に補う．
どちらの場合も $i\in J\cap K$ であり，少なくとも一方の実際の support に $i$ が属する．
補題\ref{topo:corner-singleton-radius}と包含を用いると
\[
 d_\square(x,y)
 \le d_J(x,e_i)+d_K(e_i,y)<2\sqrt5.
\]
片方の評価が真の不等式になることが，ここで必要である．
従って残るのは，両点の support がそれぞれ三元で互いに素な場合である．
座標を置換し，値 $1$ の座標をそれぞれ $1,4$ として
\[
 x=e_1+a e_2+b e_3\in Q_{1;23},\qquad
 y=e_4+c e_5+e e_6\in Q_{4;56},\qquad
 0\le a,b,c,e\le1
\]
と書く．

\emph{四通りの展開による上界．}
始点 square を平面の $[0,1]^2$ とし，終点 square を
$[-3,-2]\times[-1,0]$ に置く．三回の left step と一回の down step の
四つの順序は，次の edge-adjacent galleries として全て存在する：
\begin{align*}
 LLLD:&\quad Q_{1;23},\ Q_{1;34},\ Q_{4;13},\ Q_{4;36},\ Q_{4;56},\\
 LLDL:&\quad Q_{1;23},\ Q_{1;34},\ Q_{4;13},\ Q_{4;15},\ Q_{4;56},\\
 LDLL:&\quad Q_{1;23},\ Q_{1;34},\ Q_{1;45},\ Q_{4;15},\ Q_{4;56},\\
 DLLL:&\quad Q_{1;23},\ Q_{1;25},\ Q_{1;45},\ Q_{4;15},\ Q_{4;56}.
\end{align*}
各共有辺に沿う反転で展開すると，始点と終点の座標はそれぞれ
$(a,b)$ と $(-2-e,-c)$ となる．
この二点を結ぶ Euclidean segment が直線 $u=0,-1,-2$ と $v=0$ を横切る順序を取れば，
上の四つの gallery の少なくとも一つに segment 全体が含まれる．
格子頂点を通る場合は，二つの横切りのどちらを先に取ってもよい．
各 square 上の展開は等長なので
\begin{equation}\label{topo:eq:disjoint-square-upper}
 d_\square(x,y)^2
 \le (a+e+2)^2+(b+c)^2
 \le 4^2+2^2=20.
\end{equation}
最後の等号には $a=b=c=e=1$ が必要である．
したがって全体の直径は $2\sqrt5$ 以下で，等号点対は定理の候補に限られる．

\emph{八種類の三角形による下界．}
$d=6$，$A=\{1,2,3\}$，$B=\{4,5,6\}$ とする．
空でない $S\subset A\cup B$，$|S|\le3$ に対応する頂点に対し，
$a(S)=|S\cap A|$，$b(S)=|S\cap B|$ と置く．
各 square を singleton と triple の頂点を結ぶ対角線で二分し，
次の頂点値を各三角形上 affine に延長する：
\[
\begin{array}{c|rrrrrrrrr}
 (a,b)&(3,0)&(2,0)&(1,0)&(2,1)&(1,1)&(1,2)&(0,1)&(0,2)&(0,3)\\\hline
 \sqrt5\,h&-5&-3&-2&-1&0&1&2&3&5.
\end{array}
\]
三角形は singleton から pair，triple への subset chain であり，
加える label の $A/B$ 型は八通りである．
対応する $\sqrt5\,h$ の二つの直交辺方向の増分は
\[
\begin{array}{c|rrrrrrrr}
 \text{型}&AAA&AAB&ABA&ABB&BAA&BAB&BBA&BBB\\\hline
 \text{増分}&(-1,-2)&(-1,2)&(2,-1)&(2,1)&(-2,-1)&(-2,1)&(1,-2)&(1,2).
\end{array}
\]
従って全三角形上で $|\nabla h|^2=(1^2+2^2)/5=1$ であり，
$h$ は貼合せの長さ距離に関して $1$-Lipschitz である．
しかし
\[
 h(v_B)-h(v_A)=\frac{10}{\sqrt5}=2\sqrt5
\]
なので $d_\square(v_A,v_B)\ge2\sqrt5$ となる．
命題\ref{topo:coordinate-retraction}の等長性により，この下界は
$d>6$ の中の任意の六 label にも保たれる．
上界と座標置換対称性を合わせると，直径と全等号点対の記述を得る．
\end{proof}

\begin{corollary}[六次以上の直径対の線形安定性]
\label{topo:diameter-six-stability}
$d\ge6$，$0\le\delta<\sqrt{20}-\sqrt{13}$ とし，
$d_\square(x,y)\ge\sqrt{20}-\delta$ とする．
すると $A=\operatorname{supp}(x)$ と $B=\operatorname{supp}(y)$ は
互いに素な三元集合であり，
\[
 d_\square(x,v_A)+d_\square(y,v_B)
 \le\sqrt{20}\,\delta-\frac{\delta^2}{2}
 \le\sqrt{20}\,\delta.
\]
\end{corollary}
\begin{proof}
命題\ref{topo:coordinate-retraction}と，既存の五次の直径定理
\cite[\texttt{thm:intro-diameter-C5}, \texttt{sec:finite-diameter}]{KECubicalCurrent}
により，support の和が五元以下なら距離は $\sqrt{13}$ 以下となるので，
両 support は互いに素な三元集合である．上の座標を用い，
$u=2-a-e$，$v=2-b-c$ とすると $0\le u,v\le2$ であり，
\[
 20-\bigl((a+e+2)^2+(b+c)^2\bigr)
 =(8u-u^2)+(4v-v^2)\ge2(u+v).
\]
\eqref{topo:eq:disjoint-square-upper}から
$u+v\le(20-d_\square(x,y)^2)/2$ となる．
各点から対応する頂点までの square 内の線分の長さは，
それぞれ $2-a-b$，$2-c-e$ 以下なので，その和は $u+v$ 以下である．
距離の仮定を代入すれば結論を得る．
\end{proof}

\begin{remark}[五次の直径の使用範囲]
定理\ref{topo:diameter-six-labels}の直径と全等号点対の証明は，
上の展開・calibration と座標 retraction だけで自己完結し，五次の厳密直径を使わない．
既存の五次の $\sqrt{13}$ を使用するのは，
系\ref{topo:diameter-six-stability}の明示的な gap
$\sqrt{20}-\sqrt{13}$ の箇所に限る．
Actual KE 曲面の直径漸近を得るには，全曲面の distortion が消えることを
別に証明しなければならない．
\end{remark}

\subsection{分岐基底の立方体距離への定量的収縮}

上の $p_T$ は同相でなくても，基底上に\emph{人工的な}長さ距離を
入れると KE 側の候補立方体複体へ収束させられる．
これは収縮する方向を正確に特定する位相幾何学的な模型であり，
実際の Lagrangian fibration が誘導する action metric の収束ではない．

\begin{theorem}[有限分岐基底の立方体計量への収縮]
\label{topo:metric-collapse}
$d\ge5$ と三価木の系 $T=(T_{ij})$ を固定する．
$\mathcal C_d$ に各 unit square の Euclidean length metric
$d_\square$ を入れ，$p_T:B(T)\to\mathcal C_d$ を
系\ref{topo:KE-cubical-comparison} の全射とする．
$B(T)$ の，セルと $p_T$ に適合する有限細分上に
任意の補助 piecewise Euclidean metric $h$ を固定する．
$\varepsilon\in(0,1]$ に対して，各細分セルの接ベクトル $v$ に
\[
 |v|_\varepsilon^2
 =|d p_T(v)|_\square^2+\varepsilon^2|v|_h^2
\]
を与え，これを貼った内在的長さ距離を $d_\varepsilon$ とする．
このとき $C=C(d,T,h)<\infty$ が存在し，任意の
$x,y\in B(T)$ について
\[
 0\le d_\varepsilon(x,y)
       -d_\square(p_Tx,p_Ty)\le C\varepsilon .
\]
従って
\[
 d_{\mathrm{GH}}\bigl((B(T),d_\varepsilon),
                  (\mathcal C_d,d_\square)\bigr)
 \le \frac{C\varepsilon}{2}\longrightarrow0 .
\]
$d\ge6$ では各 $\varepsilon>0$ の分岐基底は
$\mathcal C_d$ と同相ではないにもかかわらず，この収束は成立する．
\end{theorem}

\begin{proof}
$\rho_d\circ p_T=q$ は各三角形セル $\tau_{ijk}\subset B(T)$
から $S_d$ の三角形 $[ijk]$ への homeomorphism である．
$\rho_d^{-1}$ は $x\mapsto x/\max_i x_i$ なので，
$[ijk]$ を三つの最大座標領域に分ければ smooth かつ Lipschitz である．
したがって $p_T|_{\tau_{ijk}}$ の逆写像は
$\mathcal C_d$ の対応する三つの unit squares 上で
区分的 Lipschitz となり，補助 metric $h$ に関する
一様定数 $L=L(d,T,h)$ を持つ．
$p_T^{-1}(z)$ は，$z$ が support 二の辺の内部にあるとき，
対応する $T_{ij}^{\rm int}$ の有限木（longitudinal 座標を固定した
各二角形の横断部分）であり，それ以外は一点である．
従って全ての fiber の内在的な $h$-直径は一様に
$D=D(d,T,h)$ 以下である．Fiber 内の任意の接続曲線の
$d_\varepsilon$-長は高々 $\varepsilon D$ にできる．

各 unit square は ambient $\mathbb R^d$ 内の凸集合である．
$\mathcal C_d$ の長さ最短路が同じ square を離れて再び戻るなら，
その間を square 内の線分で置き換えると短くなる
（元の任意の曲線の長さは ambient Euclidean chord 以上）．
従って最短路は，少なくとも一つの選択では，
有限個の square の gallery として記述でき，
系\ref{topo:six-label-gallery}により，その訪問数は
$N_d=3\binom{\min\{d,6\}}{3}$ 以下に取れる．
Edge に沿う部分もそれを含む任意の square に割り当てればよい．

$p_Tx$ と $p_Ty$ を結ぶこの gallery の各 square 内の線分を，
対応する $\tau_{ijk}$ の $p_T$ による逆像へ持ち上げる．
その $h$-長は元の線分長の $L$ 倍以下なので，
\[
 L_\varepsilon(\widetilde\gamma_{\rm square})
 \le L_\square(\gamma_{\rm square})
      +\varepsilon L\,L_\square(\gamma_{\rm square}).
\]
異なる square に移る点では，持ち上げた端点同士は
同じ $p_T$-fiber にある．この二点をその木の中で
$\varepsilon D$ 以下の費用で接続する．
最初と最後にも同様の fiber connector を使う．
全誤差は
\[
 \varepsilon\bigl(L\operatorname{diam}(\mathcal C_d)
                          +(N_d+1)D\bigr)
\]
以下である．逆向きの不等式は
$|dp_T(v)|_\square\le |v|_\varepsilon$ から従う．
$p_T$ は全射なので，その graph correspondence の distortionは
$C\varepsilon$ 以下となり，GH評価を得る．
\end{proof}

\begin{proposition}[収縮速度の次数による鋭い差]\label{topo:metric-collapse-sharp}
定理\ref{topo:metric-collapse}と同じ$p_T,h,d_\varepsilon$を固定し，
$p_T$に付随するgraph correspondenceのdistortionを
\[
 \mathcal D_\varepsilon=
 \sup_{x,y\in B(T)}
 \left|d_\varepsilon(x,y)-d_\square(p_Tx,p_Ty)\right|
\]
とする．$\varepsilon\downarrow0$に対し，
\[
 \mathcal D_\varepsilon=
 \begin{cases}
  \Theta_{T,h}(\varepsilon^2),&d=5,\\
  \Theta_{d,T,h}(\varepsilon),&d\ge6.
 \end{cases}
\]
従って五次では
$d_{\mathrm{GH}}((B(T),d_\varepsilon),(\mathcal C_5,d_\square))
=O_h(\varepsilon^2)$ に改善できる．
$d\ge6$ での下界は\emph{このgraph correspondenceのdistortion}に
対するものであり，別のcorrespondenceによって定義される
$d_{\mathrm{GH}}$の同じ次数の下界を主張しない．
\end{proposition}
\begin{proof}
任意の$h$-rectifiable曲線$\gamma$について積分の三角不等式を用いると，
\[
 L_\varepsilon(\gamma)
 =\int\sqrt{|(p_T\gamma)'|_\square^2
                  +\varepsilon^2|\gamma'|_h^2}
 \ge\sqrt{L_\square(p_T\gamma)^2+
                          \varepsilon^2L_h(\gamma)^2}.
\]
従って全$x,y$に対し
\begin{equation}\label{topo:metric-lower-mixed}
 d_\varepsilon(x,y)
 \ge\sqrt{d_\square(p_Tx,p_Ty)^2+\varepsilon^2d_h(x,y)^2}.
\end{equation}
$d\ge6$なら$d-4\ge2$であり，少なくとも一つの内部木
$T_{ij}^{\rm int}$ は正の長さを持つ．同じ$p_T$-fiberの相異なる
$x,y$を固定すると，\eqref{topo:metric-lower-mixed}により
$\mathcal D_\varepsilon\ge\varepsilon d_h(x,y)>0$である．
上界は定理\ref{topo:metric-collapse}による．

$d=5$では全ての内部木は一点であり，$p_T$は同相である．
有限細分上で$p_T$の逆はpiecewise Lipschitzなので，ある$L<\infty$に対し
$|v|_h\le L|dp_T(v)|_\square$が全細分セルで成り立つ．
$\mathcal C_5$内の測地線をこの同相の逆で持ち上げれば
\[
 d_\varepsilon(x,y)\le
   \sqrt{1+\varepsilon^2L^2}\,d_\square(p_Tx,p_Ty)
\]
となる．この上界と$d_\square$の直径の有限性から
$\mathcal D_\varepsilon=O(\varepsilon^2)$を得る．
また任意の相異なる$x,y$を固定すると，
\eqref{topo:metric-lower-mixed}と
$d_\square(p_Tx,p_Ty)>0,\ d_h(x,y)>0$により
$\mathcal D_\varepsilon\ge c_{x,y,h}\varepsilon^2$
が十分小さい$\varepsilon$に対して成立する．
\end{proof}

\begin{proposition}[有限squareに対するthree-input比較原理]
\label{topo:finite-portal-criterion}
$T$を固定し，$p_T:B(T)\to\mathcal C_d$を上の写像とする．
$B(T)$上の任意の長さ距離$d_t$に対し，非負の定数
$\alpha_t,\sigma_t,\lambda_t,\kappa_t$が次を満たすと仮定する．
\begin{enumerate}
\item[\textup{(N)}]
任意の$x,y\in B(T)$に対し
$d_t(x,y)\ge d_\square(p_Tx,p_Ty)-\alpha_t$
（大域no-shortcut）．
\item[\textup{(S)}]
任意のunit square $Q\subset\mathcal C_d$と任意の線分
$\gamma\subset Q$に対し，$p_T$の下で$\gamma$に一致する
rectifiable lift $\widetilde\gamma$が存在し，
\[
 L_{d_t}(\widetilde\gamma)\le
 (1+\sigma_t)L_\square(\gamma)+\lambda_t
\]
（各squareの水平持ち上げ）．
\item[\textup{(F)}]
全ての$z\in\mathcal C_d$で
$\operatorname{diam}_{d_t}p_T^{-1}(z)\le\kappa_t$
（分岐木fiber・portalの接続費用）．
\end{enumerate}
$N_d=3\binom{\min\{d,6\}}{3}$, $D_d=\operatorname{diam}(\mathcal C_d,d_\square)$
とすると
\[
 \sup_{x,y\in B(T)}
 |d_t(x,y)-d_\square(p_Tx,p_Ty)|
 \le\max\{\alpha_t,\,
       \sigma_tD_d+N_d\lambda_t+(N_d+1)\kappa_t\}.
\]
右辺が零へ行けば
$(B(T),d_t)\to(\mathcal C_d,d_\square)$ が
Gromov--Hausdorff収束する．
\end{proposition}
\begin{proof}
下界は\textup{(N)}．上界には$p_Tx$と$p_Ty$を結ぶ
系\ref{topo:six-label-gallery}の最短路を取り，$N_d$個以下の
square内線分からなるgalleryを選ぶ．
各線分を\textup{(S)}で持ち上げ，隣り合う線分の持ち上げ端点が
属する同一fiberの二点を，\textup{(F)}で評価される ambient 距離を
任意の正誤差で近似する経路で接続し，最初と最後も$x,y$へ接続する．
これらの経路が fiber 内に留まる必要はない．近似の誤差を零へ送れば，$d_t(x,y)$ は
\[
 (1+\sigma_t)d_\square(p_Tx,p_Ty)
   +N_d\lambda_t+(N_d+1)\kappa_t
\]
以下である．$p_T$は全射なので，graph correspondenceの
distortion boundからGH収束が従う．
\end{proof}

\begin{remark}[端点だけを用いる有限 gallery の上界]
\label{topo:finite-endpoint-comparison}
上界には target 線分へ厳密に写る lift を必要としない．一般の距離空間
$(X,d_X)$，写像 $F:X\to\mathcal C_d$ と非負定数 $a,\lambda,\kappa$ に対し，
各 $x$ の anchor $q_x$ を
$d_\square(F(x),q_x)\le a$ となるように選ぶ．各 unit square $Q$ と $z\in Q$ に
portal $p_Q(z)\in X$ を選び，
\[
 \begin{aligned}
 d_X(p_Q(z),p_Q(z'))&\le |z-z'|+\lambda &&(z,z'\in Q),\\
 d_X(p_Q(z),p_{Q'}(z))&\le\kappa &&(z\in Q\cap Q'),\\
 d_X(x,p_Q(q_x))&\le\kappa &&(q_x\in Q)
 \end{aligned}
\]
が成立するとする．点の選択に連続性は不要であり，一つの比較に必要な有限個の
端点に対して整合的に選べれば足りる．このとき，前命題と同じ $N_d$ に対し
\[
 d_X(x,y)\le d_\square(F(x),F(y))
       +N_d\lambda+(N_d+1)\kappa+2a .
\]
実際，$q_x,q_y$ を結ぶ $M\le N_d$ 本の最短 gallery に第一の不等式を使い，
$M-1$ 個の junction と初め・終わりに残る二つの不等式を使う．
最後に anchor の誤差 $2a$ を加えればよい．$q_x=q_y$ の場合は一つの square の
定値線分を用いる．全て ambient 距離に対する三角不等式なので，
portal 間の経路を fiber 又は局所領域の中へ閉じ込める仮定も不要である．
\end{remark}

\begin{example}[各squareの局所モデルだけでは不十分]
\label{topo:portal-counterexample}
$Q=[0,1]^2$ の二つのコピー$Q_1,Q_2$を，原点だけで貼り合わせた
長さ空間$Y=Q_1\vee_{\{0\}}Q_2$を考える．
両コピー上でidentityとなる折り畳み写像$p:Y\to Q$を取る．
この写像は非伸長なので，命題\ref{topo:finite-portal-criterion}の
\textup{(N)}は$\alpha=0$で成立する．
各squareの直線はどちらのコピーにも等長に持ち上がるので，
\textup{(S)}も$\sigma=\lambda=0$で成立する．
しかし二つの右上頂点は同じ$p$-fiberにあり，
$Y$での距離は$2\sqrt2$である．
従って \textup{(F)} は破れ，
\[
 \operatorname{diam}Y=2\sqrt2,\quad
 \operatorname{diam}Q=\sqrt2,\quad
 d_{\mathrm{GH}}(Y,Q)\ge\frac{\sqrt2}{2}.
\]
特に有限個の局所的なsquare isometriesと大域no-shortcutだけから
global GH収束は導けない．
\end{example}

\begin{example}[飽和した局所領域の外側を通る薄い shortcut]
\label{topo:protected-thin-bypass}
局所の protected distance，square の持上げ上界，小さい全 fibers だけでは，
全域 no-shortcut は従わない．この点を一般の compact length space で確認する．
Target を二枚の unit square からなる Euclidean rectangle
\[
 Y=[0,1]_s\times[-1,1]_u
\]
とし，$0<\varepsilon<1/4$ に対して source は同じ集合に
\[
 g_\varepsilon=du^2+a_\varepsilon(u)^2ds^2
\]
を入れた $X_\varepsilon$ とする．$a_\varepsilon$ は
$\varepsilon\le a_\varepsilon\le1$ を満たす滑らかな関数で，
$u\le\varepsilon/3$ と $u\ge4\varepsilon/3$ では $1$，
$2\varepsilon/3\le u\le\varepsilon$ では $\varepsilon$ に等しいものを取る．
境界と corner を許した Riemannian length metric を用い，
$f_\varepsilon:X_\varepsilon\to Y$ は集合としての恒等写像とする．

全 fibers は一点なので，その ambient 径は零である．
$g_\varepsilon\le g_Y$ より，target の各 square 内線分は恒等写像で
長さを増やさず持ち上がる．従って\textup{(S)}の誤差と\textup{(F)}の誤差は共に零である．
また
\[
 W=[0,1]\times[-1,0],\qquad
 V_\varepsilon=[0,1]\times[-1,\varepsilon/3]
\]
は $f_\varepsilon$-飽和であり，その上では metric は Euclidean なので
\[
 d_W^{\rm int}(x,y)=d_{V_\varepsilon}^{\rm int}(x,y)
 =d_Y(f_\varepsilon(x),f_\varepsilon(y))
 \qquad(x,y\in W).
\]
しかし $x=(0,0)$，$y=(1,0)$ を，上端の $u=\varepsilon$ まで上がり，
そこで水平移動して下りる経路で結ぶと
\[
 d_{X_\varepsilon}(x,y)\le2\varepsilon+\varepsilon=3\varepsilon,
 \qquad d_Y(f_\varepsilon(x),f_\varepsilon(y))=1.
\]
従って全域 no-shortcut の誤差は零へ行かない．

これは actual KE 退化の反例ではなく，局所的な protected 比較と
全域の\textup{(S),(F)}から\textup{(N)}を形式的に推論できないことを示す例である．
局所領域 $W$ 自体の飽和性も満たすが，上半分には同じ protected 比較を与えていない．
従って，飽和した\emph{全域被覆}上の局所距離比較と十分な overlap 制御を
さらに仮定した命題を否定する例とは扱わない．
本研究でも，全域被覆・全 overlap・actual 下界の確認を省略しない．
\end{example}

\begin{remark}[実際の幾何計量への接続に残るもの]
定理\ref{topo:metric-collapse}の人工計量では
\textup{(N)}が$\alpha=0$で成り立ち，
\textup{(S)}の誤差は$O(\varepsilon)$，
\textup{(F)}も$O(\varepsilon)$である．
しかしactual Lagrangian action metricやKE退化計量が
\textup{(N),(S),(F)}を満たすことは未証明である．
特に局所モデルの各squareに距離比較があるだけでは
分岐木fiber全体とportalを短く接続できるとは限らない．
逆に，$p_T$のgraph correspondenceが距離distortion
$\delta_t\to0$を持つなら，同一fiberの任意の二点$x,y$について
$d_t(x,y)\le\delta_t$なので，\textup{(F)}の$\kappa_t\to0$は
\emph{$p_T$を漸近写像に固定する限り必要}である．
例\ref{topo:portal-counterexample}は，
局所のisometric liftsと大域no-shortcutが揃っていても
この必要条件を得られないことを示す．
この一般原理の新規性は主張せず，対象研究の未解決な
幾何学的入力を明確化するために用いる．
\end{remark}

\begin{remark}[計量収縮の範囲と新規性]
これは $T$ を固定した有限 PL/CW 複体に対する直接構成で，
$\varepsilon$ は幾何から導いた parameter ではなく任意に入れた値である．
従って actual action-affine metric，KE計量，symplectic transport の
いずれについても新たな収束定理を主張しない．
この種の有限複体のcollapse自体の先行研究との新規性は未確認である．
\end{remark}

\subsection{KE 近似写像と Lagrangian fibration の適合性}

\begin{proposition}[近似写像の一様比較による fiber 収縮]
\label{topo:ke-map-compatibility}
$T$ を固定し，$p_T:B_d(T)\to\mathcal C_d$ を
系\ref{topo:KE-cubical-comparison}の全射とする．
コンパクト距離空間 $(X_t,\bar d_t)$，全射
$\pi_t:X_t\to B_d(T)$，写像 $\Phi_t:X_t\to\mathcal C_d$ が与えられ，
\[
 \sup_{x,y\in X_t}
 \bigl|\bar d_t(x,y)-d_\square(\Phi_t(x),\Phi_t(y))\bigr|\le a_t,
 \qquad
 \sup_{x\in X_t}d_\square(\Phi_t(x),p_T\pi_t(x))\le b_t
\]
を満たすとする．$f_t=p_T\circ\pi_t$ と置くと，
\[
 \operatorname{dis}(f_t)\le a_t+2b_t,\qquad
 d_{\mathrm{GH}}\bigl((X_t,\bar d_t),(\mathcal C_d,d_\square)\bigr)
 \le \frac{a_t+2b_t}{2},
\]
および
\[
 \sup_{c\in\mathcal C_d}\operatorname{diam}_{\bar d_t}f_t^{-1}(c)
 \le a_t+2b_t,\qquad
 \sup_{b\in B_d(T)}\operatorname{diam}_{\bar d_t}\pi_t^{-1}(b)
 \le a_t+2b_t
\]
が成り立つ．直径は全空間の距離を制限したものであり，
fiber 内の経路長距離ではない．特に $a_t,b_t\to0$ なら $f_t$ は GH 近似を与える．
\end{proposition}
\begin{proof}
任意の $x,y$ に対して三角不等式から
\[
 \bigl|d_\square(\Phi_t(x),\Phi_t(y))
       -d_\square(f_t(x),f_t(y))\bigr|\le2b_t
\]
なので，距離歪みは $a_t+2b_t$ 以下である．
$f_t$ は全射だから，その graph は両空間間の correspondence であり，
GH 距離の評価が従う．$f_t(x)=f_t(y)$ の場合に同じ評価を使うと
$\bar d_t(x,y)\le a_t+2b_t$ を得る．
各 $\pi_t$-fiber はある $f_t$-fiber に含まれるので最後の評価も従う．
\end{proof}

\begin{remark}[実計量への適合性と特異 fibers]
KE 退化では $\bar d_t=R_t^{-1}d_{g_t}$，
$R_t=\log(-\log|t|^2)$ を取る．
関連研究\cite{KECubicalCurrent}の $\mathrm{(CP)}_d+\mathrm{(PBE)}_d$ は
適切な $\Phi_t$ に対する $a_t\to0$ を条件付きで与えるが，
本稿の symplectic construction から $b_t\to0$ はまだ示されていない．
Moser transport の一様な距離制御も別途必要であり得る．
本命題で $a_t,b_t\to0$ が成り立てば，pinched torus，
$S^1\times\figeight$，core skeleton を含め，全 fiber の縮小を結論する．
一方，regular fibers の点ごとの縮小や，特異値集合の低次元性だけでは
この一様評価を得られない．また $d\ge6$ では，
$p_T$ の同じ木 fiber 上にある複数の $\pi_t$-fibers の集合全体も潰す必要がある．
この比較原理に一般的新規性は主張しない．
\end{remark}

\section{六次の場合の同時 Petersen 化}\label{topo:petersen-section}

$d=6$ のとき，各 $T_{ij}$ は四つの葉を二組に分けるだけで定まり，
各 $G_i$ は $K_5$ の各頂点を二個の三価頂点に展開したグラフである．
分割を独立に選ぶとリンクの同型型は一定ではないが，六つのリンクを同時に
Petersen グラフにできる．しかも，そのような分割系は完全に分類できる．

\begin{definition}
$K_6$ の完全マッチングは，六つのラベルを三組に分けたものである．
$K_6$ の one-factorization は，その十五本の辺を五個の完全マッチングに分解したものである．
\end{definition}

\begin{theorem}\label{topo:petersen}
ラベルを固定した $d=6$ の分割系で，全ての $G_i$ が Petersen グラフとなるものは
ちょうど六個である．それらは $K_6$ の六個の one-factorization と自然に対応する．
対応は，辺 $ij$ を含む完全マッチングの残りの二組を $T_{ij}$ の分割に取ることで与えられる．
ラベルの置換を許せば，この六個は一つの同型類をなす．
\end{theorem}

まず存在とその対称性を明示する．六つのラベルを
$\{\infty\}\cup\mathbb F_5$ とし，$c\in\mathbb F_5$ に対して
\[
 M_c=\bigl\{\{\infty,c\},\{c+1,c-1\},\{c+2,c-2\}\bigr\}
\]
と置く．五個の $M_c$ は one-factorization $\mathcal F$ をなす．
$i=\infty$ のリンクでは，$T_{\infty c}$ の二つの頂点を，葉 $c\pm1$ に接する
$A_c$ と，葉 $c\pm2$ に接する $B_c$ と書ける．その辺は
\[
 A_cB_c,\qquad A_cA_{c+1},\qquad B_cB_{c+2}\quad(c\in\mathbb F_5)
\]
であり，これは Petersen グラフの標準表示である．
平行移動 $c\mapsto c+a$ は $\mathcal F$ を保ち，$\infty$ を固定する．
反転 $c\mapsto c^{-1}$（$0\leftrightarrow\infty$）も $\mathcal F$ を保つ．
実際，反転は $M_0$ を固定し，$M_1\leftrightarrow M_3$，
$M_2\leftrightarrow M_4$ と入れ替える．したがって，この分割系の対称群は
六つのラベル上で推移的に作用し，全てのリンクが Petersen グラフになる．

\begin{lemma}\label{topo:factorizations}
$K_6$ の完全マッチングは十五個あり，各完全マッチングはちょうど二個の
one-factorization に属する．この二個はその完全マッチング以外を共有せず，
その中の任意の一辺の両端を入れ替える置換で互いに移る．
したがって one-factorization は六個で，互いにラベルの置換で同型である．
\end{lemma}

\begin{proof}
完全マッチングの数は $5\cdot3\cdot1=15$ である．一個を固定し，その三組を
$\{A_0,A_1\}$，$\{B_0,B_1\}$，$\{C_0,C_1\}$ と書く．
これと辺を共有しない完全マッチングは
\[
 \{A_aB_b,\ B_{1-b}C_c,\ C_{1-c}A_{1-a}\},
 \qquad(a,b,c)\in\{0,1\}^3
\]
の八個である．二個が辺を共有するのは，添字の三成分のうち相異なるものが
ちょうど一個のときである．したがって，残りを四個の互いに素な完全マッチングに
分解する問題は，三次元 cube の四頂点の独立集合を選ぶ問題になる．
そのような独立集合は偶数パリティと奇数パリティの二個だけである．
例えば $000$ を含むものは，$111$ を含めると大きさ二以下になるので，
残る三個は $011,101,110$ に限られる．
一組の両端交換は一個の添字を反転し，この二個を入れ替える．
従って one-factorization の総数は $15\cdot2/5=6$ である．
固定マッチングのラベルの置換と両端交換を使えば，全てが同型であることも分かる．
\end{proof}

\begin{lemma}\label{topo:local-petersen}
一つのラベル $i$ を固定する．$G_i$ が Petersen グラフであるための必要十分条件は，
$j\ne i$ に対する五個の完全マッチング
\[
 N_{ij}=\{ij\}\cup\{T_{ij}\text{ の二組}\}
\]
が one-factorization をなすことである．
\end{lemma}

\begin{proof}
十分性は，補題\ref{topo:factorizations}の同型性と上の明示的構成から従う．
必要性のため，$G_i$ の五本の内部辺が与える完全マッチングを除く．
Petersen グラフの任意の完全マッチングの補集合は，二個の五角形である．
この事実は標準表示で直接確認できる．完全マッチングに属する spoke の数は奇数である．
三本の場合は，残る二つの添字が外側の五角形と内側の星の両方で隣接する必要があり，不可能である．
五本の場合は明らかで，一本の場合は残りの外側と内側の四頂点のマッチングが一意に決まり，
その補集合も二個の五角形になる．
Petersen グラフには三角形がないので，完全マッチングの辺はこの五角形の弦になれず，
全て二つの五角形の間を結ぶ．内部辺を縮約すると，これらの二個の五角形は
$K_5$ の辺を互いに素に分割する．
従って，一方の五角形に沿って残る五ラベルを巡回的に付け直せば，
分割系は上の $A_c,B_c$ の表示に一致する．この表示の $N_{ij}$ は
one-factorization である．
\end{proof}

\begin{proof}[定理\ref{topo:petersen}の証明]
各リンクが Petersen であると仮定し，補題\ref{topo:local-petersen}による
頂点 $i$ の one-factorization を $\mathcal F_i$ とする．
$\mathcal F_i$ と $\mathcal F_j$ は $N_{ij}$ を共有する．
ラベル $p$ を固定し，$\mathcal F=\mathcal F_p$ と置くと，
補題\ref{topo:factorizations}より，各 $j\ne p$ に対し
\[
 \mathcal F_j=\mathcal F\quad\text{又は}\quad
 \mathcal F_j=(pj)\mathcal F
\]
である．ここで $(pj)$ は両端交換の置換である．

ある $a$ について $\mathcal F_a=(pa)\mathcal F\ne\mathcal F$ だったとする．
$b\ne p,a$ を取る．もし $\mathcal F_b=\mathcal F$ なら，
$\mathcal F_a\cap\mathcal F_b$ の唯一の完全マッチングは $pa$ を含み，
$ab$ を含めないので，$N_{ab}$ の存在に矛盾する．
残る場合は $\mathcal F_b=(pb)\mathcal F$ である．
$N_{ab}$ を $(pa)$ で引き戻したものは，$\mathcal F$ に属し $pb$ を含むため，
ある残りの三ラベル $c,e,f$ を使って
\[
 N=\{pb,ac,ef\}
\]
と書ける．一方 $N_{ab}\in(pb)\mathcal F$ から
\[
 (pb)(pa)N=\{pa,bc,ef\}\in\mathcal F
\]
も従う．これは $N$ と異なるのに辺 $ef$ を共有するので矛盾である．
したがって全ての $\mathcal F_i$ は同一である．
逆向きは既に示した．六個という数とラベル置換の下での一意性は
補題\ref{topo:factorizations}から従う．
\end{proof}

\begin{corollary}[六次の同時 Petersen 系におけるラベル対称性の破れ]
\label{topo:petersen-symmetry}
$d=6$ とし，$\mathfrak T_6^{\rm Pet}$ を全ての頂点リンクが
Petersen となるラベル付き分割系の集合とする．
$|\mathfrak T_6^{\rm Pet}|=6$ であり，
$S_6$ はラベル置換によってこの集合上に推移的に作用する．
各 $T\in\mathfrak T_6^{\rm Pet}$ の安定化部分群の位数は $120$，
指数は $6$ である．
特に全 $S_6$ で不変な単一のラベル付き Petersen 分割系は存在しない．
\end{corollary}
\begin{proof}
定理\ref{topo:petersen}より全ての系は一つのラベル置換軌道に属し，
軌道の大きさは六である．従って軌道・安定化群公式から
$|\operatorname{Stab}_{S_6}(T)|=6!/6=120$ となる．
全 $S_6$ に固定された系は軌道が一点になるはずなので存在しない．
\end{proof}

\begin{remark}
KE候補$\mathcal C_6$の座標置換は全$S_6$で対称である一方，
各Petersen分岐基底の\emph{選択した木系}は
その対称性を保たない．従ってラベル付きPetersen構成の比較に
全$S_6$-equivarianceを課すなら，単一の$T$ではなく
六個の選択の置換を同時に管理する必要がある．
これは別の分岐基底やラベルを忘れた幾何構成の存在を排除しない．
\end{remark}

この同時 Petersen 化は，六次の場合の全ての core の\emph{位相的なリンク}を統一する．
それだけで格子付きファン，integral-affine atlas，あるいは対応する
シンプレクティックな core の同定を与えるわけではない．
